%! Polynomial Functions and End Behavior — Unit 2, Lesson 2.5 — Guided Notes
\documentclass[10pt]{article}
\usepackage{precalculus-article}
\usepackage{precalculus-boxes}
\newcommand{\TallMath}[1]{$\displaystyle #1\rule[-1.4em]{0pt}{3.2em}$}

\begin{document}

\pageheader{Unit 2, Lesson 2.5}{Guided Notes}
\namedateperiod

\vspace{0.12in}

\begin{objectivebox}
By the end of this lesson, I will be able to\ldots
\begin{enumerate}[itemsep=2pt]
  \item Describe the end behavior of a polynomial function and write it using
        \blank{3cm} notation. \hfill\textit{(LO 1.6.A)}
  \item Determine the end behavior of a polynomial from its \blank{2cm} and the
        \blank{3cm} of its leading coefficient. \hfill\textit{(EK 1.6.A.3)}
\end{enumerate}
\end{objectivebox}

\vspace{0.1in}

\begin{vocabbox}
\termblanklong{End behavior}
\termblanklong{Leading term}
\termblanklong{Leading coefficient}
\termblanklong{Limit notation ($\lim_{x \to \infty}$)}
\termblanklong{Dominates}
\end{vocabbox}

\vspace{0.1in}

\begin{hookbox}
A student wraps rubber bands around a ball. The radius (in mm) grows according to
$r(t) = 2t^3 - 50t^2 + 100t$.

At $t = 1{,}000$ minutes: \quad $2t^3 = $ \blank{3cm} \quad $-50t^2 = $ \blank{3cm}

What percentage of the cubic term is the $-50t^2$ term? \blank{3cm}

Which term ``wins'' as $t$ gets very large?  \blank{4cm}
\end{hookbox}

\vspace{0.1in}

\begin{notesbox}{Leading Term Dominance}
\textbf{Key Idea:} For large $|x|$, the \blank{3cm} term of a polynomial dominates
all lower-degree terms.

\medskip
\textbf{Example.} Let $p(x) = x^3 - 100x^2$.
Complete the table by computing $p(x)$ and comparing it to $x^3$:

\smallskip
\renewcommand{\arraystretch}{1.5}
\begin{tabular}{|c|c|c|c|}
\hline
$x$ & $p(x) = x^3 - 100x^2$ & $x^3$ & $\dfrac{-100x^2}{x^3} = \dfrac{-100}{x}$ \\
\hline
$100$   & \blank{2.5cm} & \blank{2.5cm} & \blank{2.5cm} \\
$1{,}000$ & \blank{2.5cm} & \blank{2.5cm} & \blank{2.5cm} \\
\hline
\end{tabular}

\medskip
As $x \to \infty$, the ratio $-100/x \to$ \blank{1.5cm}, so $p(x)$ behaves like \blank{2cm}.

\medskip
\textbf{Conclusion:} To find end behavior, look only at the \blank{3.5cm}.
\end{notesbox}

\vspace{0.1in}

\begin{notesbox}{The Four End-Behavior Cases}
The end behavior of any polynomial is determined by its \textbf{degree} (even or odd) and the
\textbf{sign} of its leading coefficient.

\medskip
\renewcommand{\arraystretch}{1.7}
\begin{tabular}{|c|c|c|c|}
\hline
\textbf{Degree} & \textbf{Leading Coeff.} & $\lim_{x \to \infty} p(x)$ & $\lim_{x \to -\infty} p(x)$ \\
\hline
Even & Positive ($a_n > 0$) & \blank{2cm} & \blank{2cm} \\
Even & Negative ($a_n < 0$) & \blank{2cm} & \blank{2cm} \\
Odd  & Positive ($a_n > 0$) & \blank{2cm} & \blank{2cm} \\
Odd  & Negative ($a_n < 0$) & \blank{2cm} & \blank{2cm} \\
\hline
\end{tabular}

\medskip
\textbf{Memory aid:} Even-degree polynomials have both ends going the \blank{1.5cm} direction.
Odd-degree polynomials have ends going \blank{2cm} directions.

\medskip
\textbf{Practice.} State the end behavior of $p(x) = -3x^4 + 2x^2 - 7$.

\smallskip
Leading term: \blank{3cm} \quad Degree: \blank{1.5cm} \quad Sign: \blank{2cm}

\medskip
$\lim_{x \to \infty} p(x) = $ \blank{2cm} \qquad $\lim_{x \to -\infty} p(x) = $ \blank{2cm}
\end{notesbox}

\end{document}
